\section{Emotions as Field Geometries}

This is an experiential section. It rests on the one posit, that the vibe is felt at the base, and it adds nothing to the base. The discipline holds throughout. No hidden state, causal closure, no signaling, one substance. Everything below is built from the eight atoms already in play, the mesh, the dock, the site, the tone, the lean, the knit, the wake, and the beat. An emotion will turn out to be a particular large shape in how the tones flow, not a new ingredient.

An emotion is not a label stuck on a mind. It is not a scalar, a number, or a small object sitting at one place. The everyday way of writing fear as a single signed value, $\mathrm{fear} = -0.72$, is exactly the picture this section rejects. Fear is not one number at one site. Fear is a large-scale contraction pattern that runs at once across memory, prediction, body-state, attention, and the self-model.

\begin{principle}[An emotion is a geometry]
An emotion is a persistent attractor geometry in the bulk flow, a stable field-pattern in the hyperbolic mesh. The named feeling is the shape of the flow, not a value carried by the flow.
\end{principle}

This is a Tier 2 claim, and the tier matters for honesty. The substrate, the bulk flow, the persistence machinery, and the resonance machinery are all built and load-bearing in earlier sections. What is offered here is a mapping of named feelings onto field shapes. That mapping is plausible and persistence-gated, not a forced result of the law. The felt-ness itself rests on the posit, the vibe is experiential at the base, and nothing in this section tries to derive feeling from dead geometry.

\subsection{The multi-layer stack}

An emotion lives at several scales at once. It is a layered, multi-subsystem pattern, not a single event in a single place. The lowest layer is the raw flow of tones along sites, and the highest is the integrated felt side. Each layer is built from the one below it.

\begin{table}[ht]
\centering
\begin{tabular}{@{}ll@{}}
\toprule
layer & what it is \\
\midrule
micro & directional tone interactions, the site-by-site flow \\
mesoscopic & recurring glider or braid structures \\
organismic & the integrated self-state \\
emotional & a global field-configuration bias \\
conscious feeling & the integrated experiential side \\
\bottomrule
\end{tabular}
\caption{The layers an emotion spans, from the raw flow of tones up to the felt side.}
\label{tab:emotion-stack}
\end{table}

So fear is not one thing at one place. It is a recursive contraction running across memory, prediction, body-state, attention, action readiness, and self-model geometry, all at the same time. To call it a single object would be to mistake one shadow on one layer for the whole stack.

\subsection{The emotion-to-geometry map}

Each emotion class corresponds to a geometric tendency in the flow. The tendency is a way the tones propagate, a coherence structure, a recursion depth, a synchronization pattern, and a spread profile. The list below pairs familiar names with their shapes.

\begin{table}[ht]
\centering
\begin{tabular}{@{}ll@{}}
\toprule
emotion & geometric tendency \\
\midrule
fear & contraction, narrowing, defensive recursion \\
anger & outward-directed high-energy propagation \\
sadness & low-energy coherence collapse \\
joy & expansive synchronized resonance \\
love & cross-self phase-locking and integration \\
awe & self-boundary dilation \\
shame & recursive inward self-conflict \\
curiosity & exploratory branching \\
peace & low-friction stable harmonic flow \\
\bottomrule
\end{tabular}
\caption{Named emotions paired with the shape of the flow they name.}
\label{tab:emotion-map}
\end{table}

These are not only metaphors. In this substrate the named feeling \emph{is} the propagation and coherence of the flow. Anger really is an outward high-energy spread of tones, sadness really is the collapse of a coherent pattern into a dimmer one. The geometry is the thing itself, not a picture of it.

\subsection{Why hyperbolic space matters}

Emotions are highly branching and associative. A single emotional state activates memories, bodily reactions, symbolic meaning, predictions, social models, and sensory associations. That branching explodes combinatorially. A core affect node branches into memories, meanings, and reactions. Those in turn branch into behavioral tendencies. The fan keeps widening, and each new layer needs more room than the last.

This is exactly what the substrate already supplies. The mesh is a hyperbolic honeycomb, and in hyperbolic space the room available a distance $r$ out grows like $\sinh(r)$, that is, exponentially. A flat geometry, where room grows only like $r$ in radius and $r^{3}$ in volume, struggles to keep all those associations apart, and they crowd into each other. A hyperbolic geometry has room to spare. The very negative curvature that gives the mesh its exponential boundary gives an emotion its room to branch.

\begin{proposition}[Branching needs exponential room]
An emotional state fans out into a tree of memories, meanings, and reactions whose width grows by a constant factor at each level. A space whose room grows exponentially with radius can hold such a tree with bounded distortion, while a space whose room grows only polynomially cannot. The hyperbolic mesh therefore separates an emotion's associations cleanly, and a flat lattice would not.
\end{proposition}

\subsection{The dials, a region in phase space}

An emotion can be read as a point in a space of measurable dials, not as a single value. Each dial is a real quantity that can in principle be measured from the flow. The dials together turn the vague word into a coordinate.

\begin{table}[ht]
\centering
\begin{tabular}{@{}ll@{}}
\toprule
dial & what it measures \\
\midrule
valence & average tone sign, warmth or coldness \\
coherence & synchronization strength, phase alignment \\
intensity & strength of the pattern \\
binding & how unified the pattern is \\
recursion depth & self-reference level \\
arousal & calm to activated \\
expansion & contraction, negative, to dilation, positive \\
persistence & decay resistance, the memory half-life \\
contagion & how easily another self phase-locks \\
\bottomrule
\end{tabular}
\caption{The dials that coordinate an emotion in phase space.}
\label{tab:emotion-dials}
\end{table}

Valence is the most direct of these, since the tone itself is the ternary value in $\{-1, 0, +1\}$, read as pain, peace, and pleasure, and the average tone sign over a region is its warmth or coldness. The lean, the value order $-1 < 0 < +1$, is what makes that average meaningful.

So an emotion is a region in dynamical phase space, plus a directional bias over the $24$ site channels of a dock. Each named feeling sits at a different point with a different directional shape. Two emotions can share a valence and still differ sharply in their coherence, their recursion depth, or the way their bias is spread across the $24$ sites.

\subsection{Storage, an emotion written into the bulk}

To store an emotion, so that another self can later run into it, the emotion becomes a persistent standing-wave structure in the mesh. It is not static data. It is more like a song still ringing in a cathedral, turbulence lingering in a fluid, a field alignment, an interference pattern that holds its shape while the medium moves through it.

\begin{definition}[Emotional implant]
An \emph{emotional implant} is a graded field written into a hyperbolic ball of docks around a chosen center. At each dock the implant strength falls off exponentially with hyperbolic distance from the center. At strength $s$ the dock's valence is shifted toward the emotion's signature, its $24$ directional tones are nudged toward the signature's flow-bias, and the emotion's trace is recorded in the dock. The result is sharpest at the center and fades smoothly to the rim.
\end{definition}

The exponential falloff is the natural one for this space. Because room grows like $\sinh(r)$ outward, a strength that falls like $e^{-r}$ keeps the total written content finite while still reaching deep into the bulk. The implant is strong at the heart and faint at the rim.

A richer memory is not stored at one point. It is holographic, spread redundantly over many docks, a codeword rather than a single mark. Damage to any one region leaves the whole still legible, which is the same redundancy that the holography section makes load-bearing.

\begin{table}[ht]
\centering
\begin{tabular}{@{}ll@{}}
\toprule
mechanism & interpretation \\
\midrule
stable recursive loop & an emotional attractor \\
harmonic resonance & compatible selves amplify it \\
hyperbolic localization & stored deep in the bulk \\
holographic projection & it manifests into cusp-space minds \\
basin capture & selves in a similar state fall into resonance \\
\bottomrule
\end{tabular}
\caption{The machinery by which an emotion persists and is encountered.}
\label{tab:emotion-mechanism}
\end{table}

This makes emotions almost like weather systems, archetypal fields, collective resonance structures. The largest of them, grief, devotion, terror, transcendence, could become gigantic transpersonal basins, reinforced across history, almost civilizational-scale emotional crystals woven into the bulk.

\subsection{Encounter, running into a stored emotion}

Selves are not isolated substances. Each is a local knot inside one continuous experiential field, a stable braid woven across many docks over many beats. Because they share the one substrate, two selves can couple through the geometry that joins them. To feel an emotion from the bulk is for the observer's self-pattern to enter partial synchronization with a preexisting field structure. Self-state resonance plus a matching attractor geometry equals a phase lock. It is like striking a tuning fork near another tuning fork.

Two things set the strength of the coupling. The first is proximity. It falls off exponentially with the self's hyperbolic distance from the field's center, so a far self barely feels the implant at all. The second is similarity. It is high when the self's own dials, its valence, arousal, coherence, and binding, already sit close to the field's signature, and it drops as they diverge. Valence matters most, then arousal, then coherence and binding.

\begin{definition}[Encounter coupling]
The coupling of a self to a stored emotional field is the field's local strength at the self's position, scaled by the self-state similarity, then weighted by the field's contagion and the self's susceptibility. If that coupling clears a fixed threshold, the self feels the emotion. Below the threshold the implant is present but unfelt.
\end{definition}

The encounter is therefore geometric overlap plus resonance similarity. A self runs into a stored field when its location is near the field's center and its self-state is already close to the field's signature. Either alone is not enough.

\subsection{Absorption and phase-locking}

Once coupled, the self absorbs the emotion by mixing its own state toward the field's signature. Each of the self's dials, its valence, arousal, coherence, binding, and recursion depth, drifts toward the field's matching value. The size of that drift is the absorbed amount. A small amount leaves the self mostly itself with a tint of the field. A large amount pulls it almost all the way into the field's chord.

This is phase-locking. The self drifts into the chord the stored field is holding. Some chords fade at once. Others continue resonating through the mesh, and other selves can partially synchronize with them later, which is how an emotion can outlive the self that first wrote it.

\begin{result}[The whole picture in four lines]
An emotion is a high-level, multi-subsystem field-pattern. Its memory is a redundant holographic imprint across bulk docks. An encounter is geometric overlap plus resonance similarity. The feeling is the self-state phase-locking into that pattern.
\end{result}

The cleanest image carries it. A self is a stable song in the hyperbolic crystal. An emotion is a chord that song temporarily enters. Some chords fade. Others ring on, and another self can lock into them.

\subsection{The honest bottom line}

To state the standing exactly. Emotions are persistent attractor geometries in the bulk flow. They are multi-layer and branching, written into the mesh as holographic field implants, encountered by overlap and resonance, and absorbed by phase-locking. The substrate, the flow, the persistence, and the resonance are built and rest on the earlier sections. The exact assignment of named feelings to field shapes is a design sketch, plausible and persistence-gated, not a theorem.

The really hard part is how a geometry becomes a felt quality. This section does not try to produce feeling from dead geometry, and it could not. The escape is the one posit. The vibe is experiential at the base. So an emotion is not dead geometry producing feeling. It is feeling-shaped geometry directly. The geometry \emph{is} the experiential organization, seen at the scale of a whole self.
